\documentclass[11pt,a4paper]{article}
\usepackage[margin=1in]{geometry}
\usepackage{amsmath,amssymb,amsthm}
\usepackage{algorithm}
\usepackage{algpseudocode}
\usepackage{booktabs}
\usepackage{hyperref}

\theoremstyle{definition}
\newtheorem{definition}{Definition}[section]
\newtheorem{theorem}{Theorem}[section]
\newtheorem{lemma}[theorem]{Lemma}
\newtheorem{remark}{Remark}[section]

\title{\textbf{Rigorous Mathematical Specification and Attention Masks of Prefix-LM and Mixture Denoisers (UL2)}}
\author{Team Scaibu \\ \small Email: \texttt{jaydeep.9664920749@gmail.com} \\ \small Architecture and Core Engine Team}
\date{October 2026}

\begin{document}
\maketitle

\begin{abstract}
Unifying language learning paradigms requires flexibly combining causal, prefix, and span-corruption attention masks. In Prefix Language Modeling, tokens within prefix length $L_{\text{prefix}}$ attend bidirectionally, while suffix tokens are constrained by a lower-triangular causal mask.
\end{abstract}

\begin{equation}
M_{i, j} = \begin{cases} 0.0, & \text{if } j < L_{\text{prefix}} \text{ or } j \le i \\ -\infty, & \text{otherwise} \end{cases}
\end{equation}

\begin{thebibliography}{9}
\bibitem{tay2022} Y.~Tay et al., ``UL2: Unifying language learning paradigms,'' in \emph{ICLR}, 2023.
\end{thebibliography}

\end{document}
